\chapter{Wild West}

\section*{Pravidlá}
\noindent Wild West Show obsahuje 10 udalostných kariet a 8 nových postáv. Postavy sa pridajú k ostatným postavám z hry BANG!.

\textbf{Príprava:} Odlož titulnú kartu Wild West Show bokom. Zvyšné udalostné karty zamiešaj lícom nadol do samostatného balíčka (WWS balíček), na spodok pridaj titulnú kartu Wild West Show. WWS balíček umiestni do stredu stola.

\textbf{Spúšťanie udalostí:} Keď hráč zahrá Dostavník alebo Wells Fargo, vezme WWS balíček pred seba, odhalí vrchnú kartu a nahlas prečíta efekt. Efekt platí, kým niekto iný nezahrá Dostavník alebo Wells Fargo — vtedy ten hráč odhalí novú vrchnú kartu, ktorá nahradí predchádzajúcu. Predchádzajúca karta sa odstráni z hry.

\textbf{Výnimka:} Titulná karta Wild West Show sa po odhalení nenahrádza a zostáva aktívna do konca hry.

\textbf{Poznámky ku kartám:} Darling Valentine: hráči si po výmene rúk navyše doberu zvyčajné 2 karty z balíčka. Dorothy Rage: ak nútený hráč požadovanú kartu nemá, ukáže svoju ruku; ak ju má, zahrá ju akoby to bol jeho ťah, cieľ výberu určuje hráč, ktorý Dorothy Rage zahral. Miss Susanna: efekt sa nevzťahuje na hráčov preskakujúcich ťah kvôli Väzeniu. Showdown: Big Spencer môže Bang! použiť ako Vedle!; Lee Van Kliff môže na aktiváciu schopnosti odhodiť ľubovoľnú kartu. Titulná karta Wild West Show: každý hráč hrá s cieľom zostať posledný — role zostávajú rovnaké, šerif nemôže ísť do Väzenia, za zabitie banditu odmena 3 kariet stále platí, po vypadnutí šerifa hra pokračuje, víťazstvo je individuálne.

\begin{figure}[!h]
    \centering
    \begin{minipage}[t]{\minipageSize}\centering
      \includegraphics[width=\cardSize\linewidth]{wildwest/06_bavaglio.png}
      \caption[Bavaglio]{Hráči nesmú hovoriť (môžu gestikulovať, stonať, ...). Každý, kto prehovorí, stratí 1 život. Či zákaz komunikácie platí aj pre písané správy (papier, mobil) je na dohode pri stole; odporúča sa písané správy tiež zakázať.}
      \skname{Roubík}
      \altnames{Gag}
    \end{minipage}
    \hfill
    \begin{minipage}[t]{\minipageSize}\centering
      \includegraphics[width=\cardSize\linewidth]{wildwest/06_camposanto.png}
      \caption[Camposanto]{Na začiatku svojho ťahu sa všetci vyradení hráči vrátia do hry s 1 životom. Role vyradených hráčov zamiešajte a náhodne rozdeľte. Hráči sa do hry vracajú stále, aj po skončení efektu, ak sú stále v hre.}
      \skname{Hřbitov}
      \altnames{Cemetery, Bone Orchard, Cintorín}
    \end{minipage}
    \hfill
    \begin{minipage}[t]{\minipageSize}\centering
      \includegraphics[width=\cardSize\linewidth]{wildwest/06_darlingvalentine.png}
      \caption[Darling Valentine]{Na začiatku svojho ťahu odhodí každý hráč všetky karty z ruky a doberie si rovnaký počet z balíčka. Hráč na ťahu po odhodení a dobratí kariet pokračuje fázou 1.}
      \skname{Miláček Valentýn}
    \end{minipage}
    \hfill
    \begin{minipage}[t]{\minipageSize}\centering
      \includegraphics[width=\cardSize\linewidth]{wildwest/06_dorothyrage.png}
      \caption[Dorothy Rage]{Hráč na ťahu môže pomenovať kartu a vybrať hráča, ktorý ju musí zahrať, ak ju má. Musí pritom vopred opísať celú akciu vrátane cieľa – nestačí povedať len „zahraj Bang!", treba povedať napríklad „zahraj Bang! na tohto hráča". Ak túto kartu hráč nemá, ukáže svoju ruku všetkým hráčom. Ak ju má, zahrá ju v rámci svojich vlastných pravidiel a parametrov, teda podľa svojho dostrelu, svojich schopností a svojich postihov či bonusov. Za všetkých okolností sa ráta, že kartu zahral tento prinútený hráč, nie ten, kto mu to nariadil. Ak prinútený hráč prehrá Duel, stratí on život. Ak prinútený hráč zahráva Dostavník alebo Paniku!, karty si tiež ťahá on, nie hráč na ťahu.}
      \skname{Zuřivá Doroty}

    \end{minipage}
\end{figure}

\begin{figure}[!h]
    \centering
    \begin{minipage}[t]{\minipageSize}\centering
      \includegraphics[width=\cardSize\linewidth]{wildwest/06_helenazontero.png}
      \caption[Helena Zontero]{Otočte vrchnú kartu z balíčka: ak sú to srdcia alebo káry, všetky role okrem šerifa sa zamiešajú a náhodne prerozdelia.}
    \end{minipage}
    \hfill
    \begin{minipage}[t]{\minipageSize}\centering
      \includegraphics[width=\cardSize\linewidth]{wildwest/06_ladyrosadeltexas.png}
      \caption[Lady Rosa del Texas]{Počas svojho ťahu si môže každý hráč vymeniť miesto s hráčom po svojej pravici, ktorý preskočí svoj najbližší ťah. Hráči si prenesú svoje karty, postavy, role aj životy. Aby sa zabránilo nekonečným slučkám, odporúča sa obmedziť počet bezprostredne po sebe idúcich použití tohto efektu najviac na počet žijúcich hráčov.}
      \skname{Lady Rúže z Texasu}
    \end{minipage}
    \hfill
    \begin{minipage}[t]{\minipageSize}\centering
      \includegraphics[width=\cardSize\linewidth]{wildwest/06_misssusanna.png}
      \caption[Miss Susanna]{Počas svojho ťahu musí každý hráč zahrať aspoň 3 karty. Kto to neurobí, stratí 1 život. Počítajú sa aj karty zahrané ešte predtým, než Miss Susanna vstúpila do hry v tom istom ťahu. Neplatí pre hráčov, ktorí preskakujú ťah kvôli Väzeniu.}
      \skname{Madam Zuzana}
    \end{minipage}
    \hfill
    \begin{minipage}[t]{\minipageSize}\centering
      \includegraphics[width=\cardSize\linewidth]{wildwest/06_regolamentodiconti.png}
      \caption[Showdown]{Každú kartu možno hrať ako kartu Bang!. Každá karta Bang! musí byť zahraná ako karta Vedle!. Big Spencer môže použiť kartu Bang! ako kartu Vedle! a Lee Van Kliff môže na aktiváciu svojej schopnosti odhodiť ľubovoľnú z týchto dvoch kariet.}
      \altnames{Regolamento di Conti}
      \skname{Zúčtování}
    \end{minipage}
\end{figure}

\begin{figure}[!h]
    \centering
    \begin{minipage}[t]{\minipageSize}\centering
      \includegraphics[width=\cardSize\linewidth]{wildwest/06_sacagaway.png}
      \caption[Sacagaway]{Všetci hráči hrajú s odhalenými kartami na ruke, okrem rolí. Ak má dôjsť k náhodnému alebo skrytému odobratiu karty z ruky, napr. pri Panike!, Cat Balou a podobných efektoch, dotknutý hráč svoju ruku dočasne skryje, premieša a karta sa odoberie bez znalosti jej identity. Po odobratí svoju ruku znovu odhalí.\cite{Wildwest}}
    \end{minipage}
    \hfill
    \begin{minipage}[t]{\minipageSize}\centering
      \includegraphics[width=\cardSize\linewidth]{wildwest/06_wildwestshow.png}
      \caption[Wild West Show]{Cieľ každého hráča je: „Zostaň posledný v hre!". Role ostávajú rovnaké, ale cieľ každého hráča je rovnaký ako pri odpadlíkovi. Schopnosti rolí sa zachovávajú. Šerif nemôže ísť do Väzenia, za zabitie banditu sa stále získava odmena a ak šerif zabije vicešerifa, jeho postih stále platí. Hra pokračuje aj po vyradení šerifa.}
      \skname{Divoký Západ}
    \end{minipage}
    \hfill
    \begin{minipage}[t]{\minipageSize}\centering
      % \includegraphics[width=\cardSize\linewidth]{}
    \end{minipage}
    \hfill
    \begin{minipage}[t]{\minipageSize}\centering
      % \includegraphics[width=\cardSize\linewidth]{}
    \end{minipage}
\end{figure}
